\documentclass[12pt]{article}
\usepackage[T1]{fontenc}
\usepackage[utf8]{inputenc}
\usepackage{lmodern}
\usepackage{textcomp}
\usepackage{enumitem}
\usepackage{geometry}
\usepackage{graphicx}
\usepackage{amsmath, amssymb}
\geometry{margin=1in}
\begin{document}
\begin{center}
Astronomy - Graduate Level Assessment 17 \
Total Marks: 40 \
Answer all questions.
\end{center}

\section*{Multiple Choice (10 marks)}
\begin{enumerate}[label=\arabic*.]
\item What is the significance of the 1:2:4 resonance in the Jupiter's moons system?
\begin{enumerate}[label=(\alph*)]
    \item The resonance pulls Io in different directions and generates heat.
    \item It makes the orbit of Io slightly elliptical.
    \item It creates a gap with no asteriods between the orbits.
    \item It prevents formation of the ring material into other moons.
\end{enumerate}
\item Why do we see essentially the same face of the Moon at all times?
\begin{enumerate}[label=(\alph*)]
    \item because the Moon has a nearly circular orbit around the Earth
    \item because the Moon does not rotate
    \item because the other face points toward us only at new moon when we can't see the Moon
    \item because the Moon's rotational and orbital periods are equal
\end{enumerate}
\item What drives differentiation?
\begin{enumerate}[label=(\alph*)]
    \item Spontaneous emission from radioactive atoms.
    \item The minimization of gravitational potential energy.
    \item Thermally induced collisions.
    \item Plate tectonics.
\end{enumerate}
\item What are the conditions necessary for a terrestrial planet to have a strong magnetic field?
\begin{enumerate}[label=(\alph*)]
    \item fast rotation only
    \item a rocky mantle only
    \item a molten metallic core only
    \item both a molten metallic core and reasonably fast rotation
\end{enumerate}
\item Which of the following lists the ingredients of the solar nebula from highest to lowest percentage of mass in the nebula?
\begin{enumerate}[label=(\alph*)]
    \item hydrogen compounds (H$_{2}$OCH$_{4}$NH$_{3}$) light gases (H He) metals rocks
    \item hydrogen compounds (H$_{2}$OCH$_{4}$NH$_{3}$) light gases (H He) rocks metals
    \item light gases (H He) hydrogen compounds (H$_{2}$OCH$_{4}$NH$_{3}$) metals rocks
    \item light gases (H He) hydrogen compounds (H$_{2}$OCH$_{4}$NH$_{3}$) rocks metals
\end{enumerate}
\end{enumerate}

\section*{True / False (10 marks)}
\textit{Answer True or False}
\begin{enumerate}[label=\arabic*.]
\setcounter{enumi}{5}
\item True or False: The black hole at the center of the Milky Way galaxy is known as Sagittarius A*.
\item True or False: A larger object, like a baked potato, retains heat longer than a smaller object due to its greater volume-to-surface area ratio.
\item True or False: Radioactive dating of lunar samples indicates that most lunar rocks are significantly older than their terrestrial counterparts.
\item True or False: When the Moon is setting at noon, it is in the third quarter phase, positioned 90 degrees behind the Sun in the sky.
\item True or False: Jupiter and other jovian planets contain a significant fraction of non-gaseous matter, despite being called 'gas giants.'
\end{enumerate}

\section*{Long Form Response (20 marks)}
\textit{Answer each question with a clear and complete explanation.}
\begin{enumerate}[label=\arabic*.]
\setcounter{enumi}{10}
\item Discuss the characteristics and significance of the Great Red Spot on Jupiter, including its formation, longevity, and impact on our understanding of atmospheric phenomena in planetary science.
\item Discuss the implications of the distance between Earth and Altair being one million times that of the Earth to the Sun, and calculate the exact distance from Earth to Altair.
\end{enumerate}

\vfill
\noindent\textit{End of Paper}
\end{document}